\documentclass[11pt]{article}
\usepackage[T1]{fontenc}
\usepackage[utf8]{inputenc}
\usepackage{lmodern}
\usepackage[a4paper,margin=27mm]{geometry}
\usepackage{amsmath,amssymb,amsthm,mathtools}
\usepackage{booktabs,microtype}
\usepackage[hidelinks]{hyperref}
\hypersetup{pdftitle={Initial entropy-production decrease for product-profile narrow beams in the Landau-Maxwell equation},pdfauthor={Xiaojian Zeng}}
\numberwithin{equation}{section}
\newtheorem{theorem}{Theorem}[section]
\newtheorem{proposition}[theorem]{Proposition}
\newtheorem{lemma}[theorem]{Lemma}
\newcommand{\R}{\mathbb R}
\newcommand{\E}{\mathbb E}
\newcommand{\tr}{\operatorname{Tr}}
\newcommand{\diag}{\operatorname{diag}}
\newcommand{\diver}{\operatorname{div}}
\newcommand{\Dzero}{D'_{0+}}
\newcommand{\normF}[1]{\left\|#1\right\|_{\mathrm F}}
\title{Initial entropy-production decrease for\\product-profile narrow beams in the\\Landau--Maxwell equation}
\author{Xiaojian Zeng\\
\small PBC School of Finance, Tsinghua University, 100190, China\\
\small \texttt{xiaojianzeng@foxmail.com}}
\date{}
\begin{document}
\maketitle

\begin{abstract}
We study the initial entropy-production derivative for the homogeneous
Landau equation with Maxwell interaction in three velocity dimensions.
The initial densities are finite mixtures of separated beams with
transverse width \(\sqrt p\), longitudinal width \(p\), and independent
local transverse and longitudinal profiles. For every fixed normalized
configuration of positive polynomial--Gaussian profiles, we derive the
full continuum expansion \(D'_{0+}(f_p)=p^{-2}\sum_j w_j C_j+O(p^{-1})\).
The coefficient contains longitudinal score skewness. A projection
inequality and transverse moment--curvature bounds give
\(C_j\leq-45/(14a_j^2)\) and \(\sum_j w_j C_j\leq-45/14\), where
\(a_j\) is the local transverse variance. Thus the initial derivative is
negative for all sufficiently small \(p\), with a threshold depending
on the fixed configuration. The proof differentiates the collision
functional and controls mixture overlap, Gaussian tails and the initial
solution derivative. It excludes this product-profile beam mechanism for
a counterexample, while leaving universal and all-time monotonicity open.
\end{abstract}

\noindent\textbf{Keywords:} homogeneous Landau equation; Maxwell molecules;
entropy production; narrow beams; anisotropic initial data.

\section{Problem and relation to previous work}

Entropy decay does not imply decay of entropy production. For the
homogeneous Landau equation with Maxwell interaction,
Tabary~\cite[Theorem 1.1]{Tabary2026} proves monotonicity when the
directional temperatures are sufficiently well distributed, and hence
after an explicit waiting time. In his time normalization, the
three-dimensional sufficient condition is
\begin{equation}\label{eq:tabary}
 T_{\max}-1\leq c_3,\qquad c_3=\frac{20-2\sqrt{70}}{15}.
\end{equation}
Silvestre~\cite[Section 1.2]{Silvestre2026} explicitly leaves the
initial-time question for arbitrary Landau--Maxwell solutions open; his
finite-angle Boltzmann counterexamples do not transfer to this operator.
The counterexample of Junn\'e, Winter and Yolda\c{s}~\cite{Junne2026}
concerns the Landau--Coulomb kernel instead.

We consider initial densities concentrated near finitely many points on
a velocity axis. Their covariance is \(\diag(p,p,3-2p)\), so their
largest directional temperature tends to three. They lie outside
\eqref{eq:tabary} for small \(p\). Within each beam we allow a general
fixed product profile, including nonradial transverse factors and
asymmetric longitudinal scores. We derive an explicit leading
coefficient for the actual initial entropy-production derivative and
prove that it is strictly negative. The density family is an initial-data
construction, not an invariant ansatz for the evolution.

The mechanism is a competition between transverse curvature,
longitudinal curvature, and longitudinal score skewness. Projection of
the longitudinal log-curvature onto the constant and score functions
controls skewness without a symmetry assumption. Transverse moment
identities then force enough curvature to make each redistributed local
coefficient negative. A partner term cancels only after both beam-label
averages; computing a component in isolation would miss this cancellation.

The elementary information inequalities, heat-flow identity and
rotation--Gaussian representation used below are established tools; see
Villani~\cite{Villani2025} and Delarue, Menozzi and
Nualart~\cite[Proposition 2.3]{DMN2014}. Our contribution is the specific
product-beam collision coefficient, its sign bound, and its transfer to
the actual continuum initial derivative. We make no claim of exhaustive
priority over consequences implicit in earlier representations.

\section{Equation, beam family and main result}

\subsection{Collision and time conventions}
Velocity is \(v\in\R^3\), with no spatial variable or external forcing.
We use dimensionless velocity units with mass one, mean zero, and energy
three. For a density \(f\), write
\begin{equation}\label{eq:moments}
 \int f\,dv=1,\qquad \int vf\,dv=0,\qquad
 T=\int vv^{\mathsf T}f\,dv,\qquad\tr T=3.
\end{equation}
The equation and kernel are
\begin{align}
 \partial_t f&=Q(f)
 =\diver_v\int a(v-w)
 \bigl[f(w)\nabla f(v)-f(v)\nabla f(w)\bigr]\,dw,
 \label{eq:equation}\\
 a(z)&=\frac{|z|^2I_3-zz^{\mathsf T}}2.
 \label{eq:kernel}
\end{align}
All velocity integrals are over the full Euclidean space of the displayed
dimension. For matrices, \(A:B=\tr(A^{\mathsf T}B)\). The entropy and
production are
\begin{align}
 H(f)&=\int f\log f\,dv,\notag\\
 D(f)&=\frac12\iint f(v)f(w)\,
  \Delta_f(v,w)^{\mathsf T}a(v-w)\Delta_f(v,w)\,dv\,dw,
 \label{eq:D}\\
 \Delta_f(v,w)&=\nabla\log f(v)-\nabla\log f(w),\qquad
 \frac{d}{dt}H(f_t)=-D(f_t).\notag
\end{align}
Convolution of the quadratic kernel gives
\begin{equation}\label{eq:diffusion}
 A[f](v)=a(v)+\frac{3I_3-T}{2},\qquad
 Q(f)=\diver(A[f]\nabla f)+\diver(vf).
\end{equation}
Since \(\diver A=-v\), the backward generator is
\(L_t\psi=A:\nabla^2\psi-2v\cdot\nabla\psi\). It implies
\begin{equation}\label{eq:covariance}
 T'=-6(T-I_3),\qquad T(t)=I_3+e^{-6t}(T(0)-I_3).
\end{equation}
These identities agree with Tabary~\cite[Lemmas 2.1--2.2]{Tabary2026}.
For the kernel without the factor \(1/2\), the rate is \(-12\), as in
Caja-Lopez et al.~\cite[Proposition 2]{Caja2025}; thus
\(t_{\rm here}=2t_{\rm unscaled}\).

\subsection{Fixed local profiles}
Fix a finite number of distinct real centers \(\mu_j\) and positive
weights \(w_j\), with
\begin{equation}\label{eq:labels}
 \sum_jw_j=1,\qquad\sum_jw_j\mu_j=0,\qquad
 \sum_jw_j\mu_j^2=3.
\end{equation}
For each label let \(P_j\) on \(\R^2\) and \(R_j\) on \(\R\) be
centered probability densities with
\begin{equation}\label{eq:profiles}
 \operatorname{Cov}(P_j)=a_jI_2,\qquad
 \operatorname{Var}(R_j)=\lambda_j>0,\qquad
 a_j>0,\qquad\sum_jw_ja_j=1.
\end{equation}
A positive polynomial--Gaussian density means a centered nondegenerate
Gaussian multiplied by a polynomial globally bounded below by a positive
constant, and divided by its integral. Centering and covariance
conditions in~\eqref{eq:profiles} are imposed on the resulting density;
its covariance need not equal that of its base Gaussian. Neither
radiality of \(P_j\) nor parity of \(R_j\) is assumed.

Set \(\bar\lambda=\sum_jw_j\lambda_j\) and define
\begin{align}
 \kappa_p&=\left(1-\frac{2p+\bar\lambda p^2}{3}\right)^{1/2},
 \qquad m_j(p)=\kappa_p\mu_j,\label{eq:kappa}\\
 F_{j,p}(v)&=p^{-2}P_j(v_\perp/\sqrt p)
                  R_j((v_3-m_j(p))/p),\qquad
 f_p=\sum_jw_jF_{j,p}.\label{eq:beam}
\end{align}
Here \(v_\perp=(v_1,v_2)\) and \(p\) labels initial data, independently
of evolution time \(t\). We take sufficiently small \(p>0\) so that
\(\kappa_p\geq1/2\). The Jacobian of the local coordinates is \(p^2\).
Equations~\eqref{eq:labels}--\eqref{eq:kappa} give exactly
\begin{equation}\label{eq:Tbeam}
 \int vf_p\,dv=0,\qquad
 T(f_p)=\diag(p,p,3\kappa_p^2+\bar\lambda p^2)
       =\diag(p,p,3-2p).
\end{equation}

For a single product factor \(P(x)R(z)\), use the notation
\begin{align}
 g&=\nabla\log P,\quad H_P=\nabla^2\log P,\quad
 t_R=(\log R)',\quad h_R=(\log R)'',\notag\\
 J_P&=\E_P\normF{H_P}^2,\quad
 U_4=\E_P|x|^4,\quad K_2=\E_P(x\cdot g)^2,\label{eq:transmom}\\
 I_R&=\E_Rt_R^2,\quad J_R=\E_Rh_R^2,\quad
 S_3=\E_Rt_R^3,\quad \mathcal A=1+2a+U_4/4.
 \label{eq:longmom}
\end{align}
Subscripts \(j\) indicate the corresponding local factor.

\begin{theorem}[Negative initial derivative for product beams]\label{thm:main}
Assume~\eqref{eq:labels}--\eqref{eq:profiles}, with every local factor
a positive polynomial--Gaussian density as specified above. The actual
solution of~\eqref{eq:equation} starting from \(f_p\) has an initial
right entropy-production derivative, and
\begin{equation}\label{eq:asymptotic}
 \Dzero(f_p)=p^{-2}\sum_jw_jC_j+O(p^{-1})\qquad(p\downarrow0),
\end{equation}
where, with \(B_j=\mu_j^2+3\),
\begin{align}
 C_j={}&-\frac{B_j^2}{2}J_{P,j}-2\mathcal A_jJ_{R,j}
 \notag\\[-2pt]
 &+\bigl(9+5\mu_j^2-\mu_j^2K_{2,j}\bigr)I_{R,j}
       -\mu_ja_jS_{3,j}.\label{eq:C}
\end{align}
Every coefficient obeys
\begin{equation}\label{eq:sign}
 C_j\leq-\frac5{28}B_j^2J_{P,j}
       \leq-\frac{45}{14a_j^2},\qquad
 \sum_jw_jC_j\leq-\frac{45}{14}.
\end{equation}
Consequently \(\Dzero(f_p)<0\) for all sufficiently small \(p>0\).
The remainder and the threshold depend on the complete fixed
configuration. There is no assertion uniform over changing profiles,
weights, widths, center gaps, or numbers of beams.
\end{theorem}

The inequality for the functional~\eqref{eq:C} also holds for regular
positive product factors having the displayed finite moments and the
integrations by parts used in Section~\ref{sec:sign}. The stronger
polynomial--Gaussian assumption ensures the actual-mixture and
initial-solution assertions in~\eqref{eq:asymptotic}.

\section{Reconstruction from the collision functional}\label{sec:coefficient}

\subsection{Exact first variation}
Put \(q=Q(f)/f\). At any time at which the differentiation is justified,
\begin{equation}\label{eq:pairvariation}
 D'(f)=\iint f(v)f(w)
 \bigl[q(v)\Delta_f^{\mathsf T}a(v-w)\Delta_f
       +2\nabla q(v)^{\mathsf T}a(v-w)\Delta_f\bigr]\,dv\,dw.
\end{equation}
Indeed, the two density variations are equal after exchanging \(v,w\),
and the two score variations are also equal because
\(\partial_t\nabla\log f=\nabla q\). The prefactor \(1/2\) in
\eqref{eq:D} then yields~\eqref{eq:pairvariation}. From
\eqref{eq:diffusion},
\begin{equation}\label{eq:q}
 q=A[f]:\bigl(\nabla^2\log f+
                  \nabla\log f\,\nabla\log f^{\mathsf T}\bigr)+3.
\end{equation}
Thus~\eqref{eq:pairvariation} differentiates the original nonlinear
collision functional, including its partner dependence; the covariance
has not been held fixed along the evolution. Section~\ref{sec:flow}
justifies this formula at the initial time for the densities in the theorem.

\subsection{Leading local matrices and partner averages}
For the leading calculation use the coordinates and derivative scaling
\[
 v=(\sqrt p\,x,m+pz),\qquad w=(\sqrt p\,y,n+p\zeta),\qquad
 \mathsf J_p=\diag(p^{-1/2},p^{-1/2},p^{-1}).
\]
At \(p=0\), with \(m=\mu_j\), \(n=\mu_k\), the constant terms of
\(p\mathsf J_pA\mathsf J_p\) and
\(p\mathsf J_pa(v-w)\mathsf J_p\) are respectively
\begin{align}
 A_0(x;m)&=
 \begin{pmatrix}
 B I_2/2&-mx/2\\-mx^{\mathsf T}/2&1+|x|^2/2
 \end{pmatrix},\qquad B=m^2+3,\label{eq:A0}\\
 K_0(x,y;m,n)&=
 \begin{pmatrix}
 (m-n)^2I_2/2&-(m-n)(x-y)/2\\
 -(m-n)(x-y)^{\mathsf T}/2&|x-y|^2/2
 \end{pmatrix}.\label{eq:K0}
\end{align}
Define the local quantities
\begin{equation}\label{eq:QLK}
 L=\frac{\Delta P}{P}=\tr H_P+|g|^2,\quad
 K=x\cdot g,\quad d_0=1+|x|^2/2,\quad
 \ell=\frac{R''}{R}=h_R+t_R^2.
\end{equation}
With \(s=(g,t_R)\), the leading term of \(pq\) is
\begin{equation}\label{eq:Q0}
 Q_0=\frac B2 L-mKt_R+d_0\ell.
\end{equation}
In this subsection \(\E_*\) averages the partner profile and then all
partner labels with their weights. Set \(\bar I=\sum_kw_kI_{R,k}\).
Centering, covariance and score integration by parts give
\begin{align}
 M_2&:=\E_*[K_0(s-s_k)]=A_0s+m e_3,\label{eq:M2}\\
 M_1&:=\E_*[(s-s_k)^{\mathsf T}K_0(s-s_k)]\notag\\
 &=\frac B2|g|^2-mKt_R+2mt_R+d_0t_R^2
                  +\frac{|x|^2}{2}\bar I+c_m.\label{eq:M1}
\end{align}
The constant in the last line is
\[
 c_m=\frac12\sum_kw_k\bigl[(m-\mu_k)^2\E_{P_k}|g_k|^2
                                   +2a_kI_{R,k}\bigr].
\]
It has no local-coordinate dependence. The score identities used here
are \(\E_{P_k}g_k=0\),
\(\E_{P_k}[y_i(g_k)_r]=-\delta_{ir}\), and
\(\E_{R_k}t_{R,k}=0\); local product independence is essential.
Since \(\diver A_0=-me_3\),
\(\diver M_2+s\cdot M_2=Q_0\). Integration of the gradient term in
\eqref{eq:pairvariation} therefore reduces the leading row to
\begin{equation}\label{eq:row}
 \mathcal R_j=\E_{P_jR_j}\bigl[Q_0M_1-2Q_0^2\bigr].
\end{equation}
Also \(\E Q_0=0\), so \(c_m\) drops out.

\subsection{Local integrations and the label cancellation}
Suppress the label again. The longitudinal identities are
\begin{align}
 &\E t_R=0,\quad \E h_R=-I_R,\quad
 \E(h_Rt_R)=-S_3/2,\quad
 \E(h_Rt_R^2)=-\E t_R^4/3,\notag\\
 &\E\ell=0,\quad \E(t_R\ell)=S_3/2,\quad
 \E(t_R^2\ell)=2\E t_R^4/3,\quad
 \E\ell^2=J_R+\E t_R^4/3.\label{eq:longibp}
\end{align}
For example, integrate the derivatives of \(Rt_R\), \(Rt_R^2\) and
\(Rt_R^3\). In the transverse variables,
\begin{align}
 &\E K=-2,\quad\E(Kd_0)=-2-4a,\quad
 \E d_0=1+a,\quad \E d_0^2=\mathcal A,\notag\\
 &\E L=0,\quad\E(Ld_0)=2,\quad\E(L|x|^2)=4,\notag\\
 &\E(L|g|^2)-2\E L^2=-2J_P.\label{eq:transibp}
\end{align}
The last identity is the usual heat-flow Fisher dissipation identity:
for \(\partial_\tau P=\Delta P\), differentiating \(\E|g|^2\) gives
\(\E(L|g|^2)+2\E(g\cdot\nabla L)\), with
\(\E(g\cdot\nabla L)=-\E L^2\); integration by parts also gives
the derivative as \(-2J_P\). The polynomial--Gaussian factors justify
both calculations directly, without log-concavity.

Substitution of~\eqref{eq:longibp}--\eqref{eq:transibp}
into~\eqref{eq:row} cancels all fourth-score moments and yields
\begin{equation}\label{eq:rawrow}
 \mathcal R_j=-\frac{B_j^2}{2}J_{P,j}-2\mathcal A_jJ_{R,j}
 +(3+5\mu_j^2-\mu_j^2K_{2,j})I_{R,j}
 -\mu_ja_jS_{3,j}+B_j\bar I.
\end{equation}
The difference between this row and \(C_j\) in~\eqref{eq:C} is
\(B_j\bar I-6I_{R,j}\). It vanishes after the final label sum:
\begin{equation}\label{eq:redistribution}
 \sum_jw_j( B_j\bar I-6I_{R,j})
   =6\bar I-6\bar I=0,
\end{equation}
because \(\sum_jw_jB_j=6\). This establishes the asserted leading
coefficient before passing from the separated component calculation
to the actual mixture.

\section{A negative bound for every local coefficient}\label{sec:sign}

\begin{lemma}[Longitudinal projection]\label{lem:long}
For a regular positive probability density \(R\) with the moments and
integrations by parts above,
\begin{equation}\label{eq:projection}
 J_R\geq I_R^2+\frac{S_3^2}{4I_R}.
\end{equation}
\end{lemma}
\begin{proof}
An integrable positive density has \(I_R>0\). In \(L^2(R\,dz)\), the
functions \(1\) and \(t_R\) are orthogonal, with squared norms \(1\)
and \(I_R\). Their inner products with \(h_R\) are \(-I_R\) and
\(-S_3/2\) by~\eqref{eq:longibp}. The squared norm of the orthogonal
projection of \(h_R\) gives~\eqref{eq:projection}.
\end{proof}

For any \(\mathcal A>0\), completing the skewness square now gives
\begin{align}
 -2\mathcal A J_R-maS_3
 &\leq-2\mathcal A I_R^2
   -\frac{\mathcal A}{2I_R}
       \left(S_3+\frac{maI_R}{\mathcal A}\right)^2
   +\frac{m^2a^2}{2\mathcal A}I_R\notag\\
 &\leq-2\mathcal A I_R^2+\frac{m^2a^2}{2\mathcal A}I_R.
 \label{eq:square}
\end{align}
This estimate retains both signs and arbitrary magnitudes of score
skewness; no symmetry of \(R\) has been introduced.

\begin{lemma}[Transverse curvature and fourth moment]\label{lem:trans}
Let \(P\) be a regular positive centered density on \(\R^2\) with
covariance \(aI_2\) and the stated finite moments and integrations by
parts. Then \(u=U_4/a^2>4\), and
\begin{equation}\label{eq:threebounds}
 J_P\geq\frac2{a^2},\qquad
 K_2\geq4+\frac{16}{u-4},\qquad
 K_2\leq4+\sqrt{2J_PU_4}.
\end{equation}
In particular,
\begin{equation}\label{eq:14}
 J_PU_4\geq14,\qquad J_P\mathcal A\geq7/2.
\end{equation}
\end{lemma}
\begin{proof}
Since \(\E|x|^2=2a\) and \(|x|^2\) is not constant almost surely
under a positive density, \(U_4>4a^2\). Integration by parts gives
\(\E(x\cdot g)=-2\); hence Cauchy--Schwarz implies
\(\E|g|^2\geq2/a\). Also
\(\E H_P=-\E(gg^{\mathsf T})\). Jensen and the matrix trace inequality
give
\[
 J_P\geq\normF{\E H_P}^2
       \geq\tfrac12(\E|g|^2)^2\geq2/a^2.
\]
Writing \(r^2=|x|^2\), another integration by parts gives
\(\E(Kr^2)=-8a\). Thus
\[
 \operatorname{Cov}(K,r^2)=-4a,\quad
 \operatorname{Var}(K)=K_2-4,\quad
 \operatorname{Var}(r^2)=U_4-4a^2.
\]
Their covariance inequality proves the lower bound on \(K_2\).
Pointwise \(K^2\leq r^2|g|^2\), while
\[
 \E(r^2|g|^2)=4-\E(r^2\tr H_P)
       \leq4+\sqrt{2J_PU_4}.
\]
This proves the upper bound. Combining them gives
\[
 J_PU_4\geq\max\left\{2u,\frac{128}{(u-4)^2}\right\}\geq14.
\]
For \(u\geq7\) use \(2u\geq14\); for \(4<u\leq7\) use
\(128/(u-4)^2\geq128/9>14\). Finally
\(\mathcal A\geq U_4/4\) proves the last assertion.
\end{proof}

\begin{proof}[Proof of the coefficient bound in Theorem~\ref{thm:main}]
Fix one label and write \(m=\mu_j\), \(B=m^2+3\). Since
\(a^2/(2\mathcal A)\leq2/u\), Lemma~\ref{lem:trans} implies
\begin{equation}\label{eq:cbound}
 5-K_2+\frac{a^2}{2\mathcal A}
 \leq1-\frac{16}{u-4}+\frac2u<1.
\end{equation}
Equations~\eqref{eq:C} and~\eqref{eq:square} consequently give
\[
 C_j\leq-\tfrac12B^2J_P-2\mathcal A I_R^2+cI_R,
 \qquad c\leq9+m^2\leq3B.
\]
As \(I_R\geq0\), maximization of the last quadratic and
\(J_P\mathcal A\geq7/2\) yield
\begin{align}
 C_j&\leq-\frac12B^2J_P+\frac{9B^2}{8\mathcal A}
       \leq-\frac5{28}B^2J_P\notag\\
 &\leq-\frac5{14}\frac{B^2}{a^2}
       \leq-\frac{45}{14a^2}.\label{eq:localbound}
\end{align}
Jensen's inequality for \(a\mapsto a^{-2}\), together with
\(\sum_jw_ja_j=1\), gives \(\sum_jw_ja_j^{-2}\geq1\) and
completes~\eqref{eq:sign}. The numerical constant is a convenient
nonsharp bound; no optimality is asserted.
\end{proof}

\section{Passage from component coefficients to the mixture}\label{sec:transfer}

\subsection{Polynomial scaling of the component calculation}
In a component-pair integral use the true whole-mixture diffusion
matrix \(A=a(v)+(3I_3-T(f_p))/2\), but replace the logarithmic
derivatives at \(v,w\) by those of \(F_{j,p},F_{k,p}\), respectively.
Let \(\mathcal P_{jk}(p)\) denote this replacement in
\eqref{eq:pairvariation}, including the corresponding replacements
in \(q\) and \(\nabla q\).

For temporarily fixed centers \(m,n\), both full scaled matrices
\(p\mathsf J_pA\mathsf J_p\) and
\(p\mathsf J_pa(v-w)\mathsf J_p\) are polynomials in \(p\) of degree
at most two. The local scores and their derivatives are independent of
\(p\). Formula~\eqref{eq:q} shows that \(pq\) is also a polynomial
of degree at most two, including the term \(3p\). In the gradient
term, \(\nabla_v=\mathsf J_p\nabla_{x,z}\); hence \(p^2\) times
the full replaced integrand has degree at most four. Every coefficient
is integrable under the fixed polynomial--Gaussian product factors.
Replacing \(m,n\) by \(\kappa_p\mu_j,\kappa_p\mu_k\), with
\(\kappa_p=1+O(p)\), changes the leading coefficient by \(O(p)\)
and leaves the other coefficient integrals bounded. Therefore
\begin{equation}\label{eq:componentexpansion}
 \sum_{j,k}w_jw_k\mathcal P_{jk}(p)
      =p^{-2}\sum_jw_jC_j+O(p^{-1}).
\end{equation}
The normalization identities for the limiting centers are used at
\(p=0\), not imposed on the finite-\(p\) centers separately from
the internal variances.

\subsection{Posterior overlap and full-space tails}
We give the bounds that permit the score replacements just made.
Let \(d_* =\min_{j\ne k}|\mu_j-\mu_k|>0\). Let
\(\sigma_{\min},\sigma_{\max}\) be the smallest and largest base
Gaussian standard deviations of the longitudinal factors, and let
\(b_{\max}\) bound all transverse base Gaussian covariance eigenvalues.
All these are fixed positive constants. Choose
\(r_0=d_*\sigma_{\min}/(16\sigma_{\max})\), and use the physical
core
\begin{equation}\label{eq:cores}
 \Omega_j(p)=\{\,|v_3-m_j(p)|\leq r_0,\quad |v_\perp|\leq1\,\}.
\end{equation}
For sufficiently small \(p\), every other center is separated from
this core in the longitudinal coordinate. Comparison of the Gaussian
exponents, allowing for unequal base variances, gives on \(\Omega_j(p)\)
\[
 \frac{\text{base Gaussian of }F_{k,p}(v)}
      {\text{base Gaussian of }F_{j,p}(v)}
       \leq C\exp(C/p-c/p^2),\qquad k\ne j.
\]
The fixed choice of \(r_0\) makes the coefficient of \(p^{-2}\)
strictly negative. The transverse ratio contributes at most
\(\exp(C/p)\). Positive lower bounds on the defining polynomials
allow their ratio to contribute only a polynomial in \(p^{-1}\) and
\(|v|\). Thus the wrong-component posterior weights are exponentially
small on each core. Raw derivative ratios through order three, and hence
the logarithmic derivatives entering \(q,\nabla q\), have polynomial
bounds in \(p^{-1}\) and \(|v|\). Differentiation of posterior weights
preserves the exponential estimate, with an additional polynomial factor.

Under component \(j\), the complement of~\eqref{eq:cores} has a
polynomial-weighted Gaussian tail bounded by a power of \(p^{-1}\)
times
\begin{equation}\label{eq:tails}
 C\exp\!\left(-\frac{r_0^2}{4\sigma_{\max}^2p^2}\right)
       +C\exp\!\left(-\frac1{4b_{\max}p}\right).
\end{equation}
One half of each Gaussian exponent supplies the displayed tail and
the other supplies any required polynomial moment. For the mixture,
\(\partial^\alpha f_p/f_p\) is a posterior-weighted average of the
component raw derivative ratios. Consequently the exact and replaced
pair integrands have global polynomial bounds of the same type, without
requiring a uniform positive lower bound on \(f_p\) over velocity space.
Split each component-pair integral into its two cores and their
complement. For fixed configuration-dependent \(C,K,c>0\), this gives
\begin{equation}\label{eq:overlap}
 \left|\mathcal V(f_p)-\sum_{j,k}w_jw_k\mathcal P_{jk}(p)\right|
       \leq Cp^{-K}e^{-c/p}=o(p^{-1}),
\end{equation}
where \(\mathcal V(f)\) denotes the right-hand side of
\eqref{eq:pairvariation}. This is an analytic full-space estimate,
not a velocity-grid or finite-difference approximation.

\section{The actual initial right derivative}\label{sec:flow}

It remains to identify \(\mathcal V(f_p)\) with the actual solution
derivative at each fixed admissible \(p\). All constants in this section
may depend on \(p\); no uniform ellipticity as \(p\downarrow0\) is used.
Put
\[
 C(t)=\frac{3I_3-T(t)}2,\qquad
 B_{ir}=\frac{v_i\partial_r-v_r\partial_i}{\sqrt2}\quad(i<r).
\]
For the covariance~\eqref{eq:Tbeam}, \(C(t)\) is positive definite
near zero. The backward generator in~\eqref{eq:diffusion} satisfies
\begin{equation}\label{eq:angular}
 L_t=A:\nabla^2-2v\cdot\nabla
       =\sum_{i<r}B_{ir}^2+C(t):\nabla^2-v\cdot\nabla.
\end{equation}
The angular diffusion has an orthogonal Brownian flow \(O_t\),
normalized to generate \(\sum B_{ir}^2\). Variation of constants
gives a process
\begin{equation}\label{eq:affine}
 V_t=e^{-t}O_tV_0+G_t,
\end{equation}
where, conditional on the angular path, \(G_t\) is a centered Gaussian
independent of \(V_0\), with covariance
\begin{equation}\label{eq:noise}
 \operatorname{Cov}(G_t\mid O)
 =2\int_0^t e^{-2(t-s)}O_tO_s^{\mathsf T}C(s)O_sO_t^{\mathsf T}\,ds.
\end{equation}
This is the established rotation--Gaussian mechanism of
Delarue, Menozzi and Nualart~\cite[Proposition 2.3]{DMN2014}, here
reconstructed with the kernel and time convention~\eqref{eq:kernel}.

For completeness this prescribed-covariance process is self-consistent.
Its second moment \(M(t)\) obeys
\[
 M'=(\tr M+3)I_3-5M-T(t),\qquad
 (\tr M)'=6-2\tr M.
\]
Thus \(\tr M=3\), and \((M-T)'=-5(M-T)\) with zero initial value
gives \(M=T\). Its density therefore solves the nonlinear equation
\eqref{eq:equation}. Conversely, a regular normalized solution has the
covariance~\eqref{eq:covariance} and hence this same linear generator;
the affine stochastic equation has a unique law.

The density formula is
\begin{equation}\label{eq:densityflow}
 f_t(v)=e^{3t}\E\bigl[f_p(e^tO_t^{\mathsf T}(v-G_t))\bigr].
\end{equation}
For fixed \(p\), every raw derivative ratio
\(\partial^\alpha f_p/f_p\), \(|\alpha|\leq4\), grows at most
polynomially, and the initial density has moments of every order.
Differentiate~\eqref{eq:densityflow}. The derivative of the affine
argument is a bounded linear map for small time. Jensen's inequality
with the posterior density weights, followed by change of variables,
therefore bounds, uniformly on a sufficiently small time interval,
\begin{equation}\label{eq:weighted}
 \int f_t(v)(1+|v|)^k
       \left|\frac{\partial^\alpha f_t(v)}{f_t(v)}\right|^s\,dv,
 \qquad |\alpha|\leq4,\quad k\geq0,\quad s\geq1.
\end{equation}
The bounds reduce to initial polynomial-weighted derivative moments
and bounded moments of the conditional Gaussian noise. Orthogonality
controls the angular factor. Extra velocity powers in~\eqref{eq:weighted}
give uniform integrability at infinity.

As \(t\downarrow0\), \(O_t\to I_3\) and \(G_t\to0\), so density
derivatives converge locally to the initial derivatives. Positivity
gives local convergence of their ratios. Compact/tail splitting and
H\"older's inequality, using~\eqref{eq:weighted}, now give convergence
of both \(D(f_t)\) and \(\mathcal V(f_t)\). At positive times, cutoff
differentiation and integration by parts prove
\(dD(f_t)/dt=\mathcal V(f_t)\); the same bounds remove the cutoffs.
Integration from \(\varepsilon\) to \(t\) and passage
\(\varepsilon\downarrow0\) yield
\[
 D(f_t)-D(f_p)=\int_0^t\mathcal V(f_s)\,ds,
 \qquad \Dzero(f_p)=\mathcal V(f_p).
\]
Together with~\eqref{eq:componentexpansion}--\eqref{eq:overlap}, this
proves~\eqref{eq:asymptotic}. The bound~\eqref{eq:sign} then completes
Theorem~\ref{thm:main}.

\section{Interpretation and limitations}

Separated counterstreaming beams with strong temperature anisotropy are
a natural test of initial entropy-production monotonicity. The theorem
shows that arbitrary fixed non-Gaussian product profiles do not produce
a positive initial derivative in the scaling~\eqref{eq:beam}, even when
longitudinal score skewness is present. For this beam geometry, a positive
sequence as \(p\downarrow0\) would have to leave at least one of the
fixed-configuration or local-independence hypotheses. These are possible
escape routes, not assertions that a counterexample exists there.

The conclusion concerns only the initial right derivative. It does not
establish monotonicity at later times, an invariant mixture family,
arbitrary correlated local profiles, rough initial data, or a statement
uniform over varying configurations. The remainder constant and finite
threshold are not computed explicitly, and \(45/14\) is not sharp.
The functional sign estimate holds under broader integration-by-parts
conditions than the actual-mixture theorem; we do not extend the latter
without its tail and regularity assumptions. No Boltzmann or Coulomb
conclusion follows by changing the kernel in the present proof. The
universal Landau--Maxwell initial-time monotonicity problem remains
unresolved by this result.

\section*{Reproducibility}
The accompanying original Python/SymPy programs reconstruct the leading
two-particle coefficient from explicit product-profile moment identities,
check the partner redistribution and square completion, and verify the
degree bounds of the full scaled matrices. The supplement specifies the
environment and one command for replay. Exact algebra checks support the
displayed calculation; the integrations by parts, mixture estimates and
solution differentiability are justified in the text. No numerical scan
is used as a proof of the sign, and no experimental data are used.

\section*{Use of generative AI tools}
OpenAI Codex was used in developing the analytic argument, generating
symbolic checking code, performing literature searches and internal
checks, and drafting this article. The disclosure applies throughout
the argument and to the accompanying code. AI checks are not external
peer review and can share errors. Human scientific review, confirmation
of authorship and affiliation, and acceptance of responsibility for the
complete manuscript remain pending for this local draft.

\bibliographystyle{plain}
\bibliography{manuscript/references}
\end{document}
